%%%%%%%%%%%%%%%%%%%%%%%%%%%%%%%%%%%%%%%%%%%%%%%%%%%%%%%%%%%%%%%%%%%%%%%%%%%%%%%
% OpenReview-ready responses (UNDER 5000 chars per reviewer)
% This file is a snapshot of the concise version.
%%%%%%%%%%%%%%%%%%%%%%%%%%%%%%%%%%%%%%%%%%%%%%%%%%%%%%%%%%%%%%%%%%%%%%%%%%%%%%%
\documentclass[11pt,a4paper]{article}

\usepackage[margin=1in]{geometry}
\usepackage[T1]{fontenc}
\usepackage[utf8]{inputenc}
\usepackage{microtype}
\usepackage{amsmath,amssymb}
\usepackage{mathrsfs}
\usepackage{enumitem}
\setlist[itemize]{nosep,leftmargin=*}

\newcommand{\NP}{Nguyen \& Pham (2023)}

\title{Author response (condensed, per-reviewer)}
\author{Submission 25189}
\date{\today}

\begin{document}
\maketitle

\begin{abstract}
\noindent
We thank the reviewers for their careful reading and constructive feedback. Below we reply \textbf{per reviewer}, following the OpenReview structure (``Comment:'' / ``Response:''), and we indicate concrete camera-ready edits. Each reviewer block is written to fit within typical OpenReview character limits for an author response.
\end{abstract}

\section*{Author response --- Reviewer VCqo}

\textbf{Comment:}
\begin{verbatim}
Some of the figures could be cleaner and clearer. As an example of this, in Figure 2, it is not clear. Why not keep a two color code. Show the high rank curve in one color (e.g. blue) and the low rank curves in another (or shades of another) color. Also if you have a low rank and a high rank curve at the same momentum, wouldn't it make sense to only keep those two curves? Does it make sense to compare curves with different momentums ?

Section 2,

you should keep the same citation style throughout the paper. Before line 126 you cite works by mentioning the authors and then providing the reference in the bibliography but then on line 130 (left column), you only provide the citation in the bibliography
line 139, frozen random features ensure supp(W^0(C_1)) —> this is not clear
line 110, second column, what do you mean by the mean-field width limit ? Do you mean the limit of infinite width networks?
Section 3.1.

I understand that the low rank factorization appears from the stacking of W_j^{(\ell)} and h^{(\ell-1)} but wouldn’t it be clearer to write x’ = \varphi(W^T A x + b) ?
Section 3.2.

Line 166, your notations with the lower and uppercase case c_1 in W^0(c_1) is not really clear. Why put that C_i ? Why not just keep the notation of Eq. 1, W_0^{(i)} ?
Your expectation is taken with respect to the neurons indices ? which is strange to me ? You also say that in the infinite width limit, the neurons indices are continuous. Isn’t the norm to work with a measure on the weight vectors associated to each neuron ? I.e whose support would be non zero at the points W(C_i) ?
Section 3.3.

lines 187-188, I strongly recommend writing down the expression of the loss at least and recall that the system corresponds to a simple gradient descent on this loss? Because the way you introduce (2) is a little rough
lines 187-188: What are the learning rate schedules?
Section 3.4.

On lines 166 - 168, you say that the “Better than null” assumption requires the initial loss to be better than the null function.. What does that mean ?
Line 178, second column, you say that for Relu, the same holds with high probability in r. I would recall the meaning of r here
I don’t know how easily this can be done but I would recommend laying out the assumptions cleary on lines 215 - 216. Some of them are not very clear. E.g. what do you mean by bounded activation and MIXING (i.e the mixing part)?
Section 3.5 - 3.6

lines 207 - 210, second column, you have |H_\ell(W’) - H_\ell(W’’)|< K||W^{\ell}||_{\infty, 1}. What is W^{\ell} in this setting? Does that mean all the weights from layer ell ? but then how does this relate to W’ and W’’? That notation does not make sense to me.
lines 218 - 219, what is the solution operator F?
lines 220 - 223: “After defining norms, the solution operator F,… and the contraction operator …” —> “After defining norms, F and bounds, and the contraction, the argument proceeds… ” (I think you want to add a coma here)
Section 3.7

In the statement of Theorem 3.2., when you talk about the mean field limit, you are referring to n_1, n_2 —> infinity right? A depth 2 network as in Theorem 3.3. Or does it hold for an arbitrary number of layers?
lines 269 -, What you call “High level proof idea” is pretty opaque. I think I would remove it. Or I would remove lines 227 - 243, second column and expand a little bit on the proof.
lines 272 - 273, “this yields \mathbb{E}_Z[upstream\times local]” —> what does that mean? I understand you are referring to the expressions in (2) but what do the terms upstream and local refer to?
lines 220 - 223: What does \partial_2 \mathcal{L} mean here? Does that mean the partial derivative with respect to the second argument of \mathcal{L}?
lines 239 - 241, second column, it is not clear to me why the frozen random features yield heterogeneous shifts across neurons. In fact, even before this, what do you mean by “shifts” in this setting?
Section 3.8

In the statement of Theorem 3.3. Here as well, it is not clear why/how you initialize over the indices.. you have a tuple {n1, n2} that is in the index set of Init ? What does that mean concretely ? Why just a two tuple? What do n1 and n2 represent ? do they refer to the number of neurons in the first and second layers ? Does the result only hold for depth 2 then?
Same Theorem, what do you mean by the “coupling procedure”?
Same, What does 
 refer to here ? If I’m not wrong, you don’t introduce it. I understand it measures the deviation between the mean field solution and the finite solution but what measure do you use?
\end{verbatim}
\textbf{Response:} We agree and will revise the camera-ready accordingly.
\begin{itemize}
  \item \textbf{Presentation.} Redraw Fig.~2 with a two-color scheme and matched-momentum comparisons; unify citation style; fix minor typos/phrasing.
  \item \textbf{Definitions.} Add explicit loss $\mathscr{L}$/$\mathcal{L}$, define $\partial_2\mathcal{L}$, define $F$ (Picard limit / flow map), and clarify the ODE--SGD discretization and learning-rate schedules $\xi(t)$.
  \item \textbf{Frozen features.} Restate the frozen-feature richness assumption in standard terms (random-feature density) and remove ambiguous ``$\operatorname{supp}$'' wording.
  \item \textbf{Assumptions.} Define ``better than null'', define ``mixing'' ($L^{(\ell)}$, $\|L^{(\ell)}\|_{\infty,1}\le rK$), and clarify ReLU vs.\ the nonzero-derivative assumption (GeLU satisfies; ReLU has a high-probability-in-$r$ heuristic under symmetric mixing).
  \item \textbf{Lipschitz bound (explicit).} In our architecture,
  \[
  H_\ell(c_\ell;x,W)=\sum_{k=1}^r L^{(\ell)}_{c_\ell,k}\,f_{\ell,k}(x;W),
  \qquad
  \|L^{(\ell)}\|_{\infty,1}:=\max_{c_\ell}\sum_{k=1}^r|L^{(\ell)}_{c_\ell,k}|.
  \]
  Hence for two trainable parameter tuples $W',W''$ (same frozen $L$),
  \[
  |H_\ell(c_\ell;x,W')-H_\ell(c_\ell;x,W'')|
  \le \|L^{(\ell)}\|_{\infty,1}\max_{1\le k\le r}|f_{\ell,k}(x;W')-f_{\ell,k}(x;W'')|.
  \]
  We will replace the ambiguous $\|W^{(\ell)}\|_{\infty,1}$ by $\|L^{(\ell)}\|_{\infty,1}$ and standardize $W$ to denote \emph{trainable} weights only.
  \item \textbf{Solution operator $F$ (explicit).} We use $F$ for the Picard (integral) map of the mean-field ODE: if $\dot\mu=\Phi(\mu)$, then
  \[
  (F\mu)(t)=\mu(0)+\int_0^t\Phi(\mu(s))\,ds.
  \]
  \item \textbf{Theorems / proof sketches.} Clarify that Theorem~3.2 is for arbitrary depth in our architecture and specify which widths diverge; shorten/move the high-level proof sketch; replace ``upstream$\times$local'' by ``backward$\times$forward''; define ``shifts'' (neuron-independent translation under i.i.d.\ collapse in \NP{}, Theorem~54 / Corollary~30) and explain heterogeneity via frozen mixing (since $H_2(c_2;x,W)=\sum_k L_{c_2,k}f_k(x;W)$ depends on $c_2$; analogous in spirit to bidirectional diversity in \NP{}, Theorem~38); rewrite Theorem~3.3 initialization language and define $D$ (Wasserstein-type distance as in \NP{}).
\end{itemize}

\section*{Author response --- Reviewer RZVa}

\textbf{Comment:}
\begin{verbatim}
While symmetry is preserved at
, it becomes largely imperceptible at
. Since
 remains an extremely low-rank constraint relative to a width of
, this sensitivity challenges the practical robustness of the feature learning mechanism .

Theorem 3.3 bounds the finite-width approximation error by
. Does the symmetry loss at
 occur because the exponential factor
 dominates
, causing divergence from idealized mean-field dynamics? If so, is "symmetry preservation" a genuine benefit of low-rank architectures, or merely a fragile artifact of keeping
 artificially small to prevent mean-field breakdown?

If
 must be kept strictly minimal to preserve mean-field behaviors, how does this framework scale to datasets with higher intrinsic dimensionality
? If evaluated on a synthetic dataset with controlled
, is there a theoretical or empirical phase transition when
? It would be insightful to clarify if the network loses its theoretical structural benefits once the required rank
 exceeds what the practical width
 can exponentially suppress.
\end{verbatim}
\textbf{Response:} We thank the reviewer for the detailed and actionable feedback.
\begin{itemize}
  \item \textbf{Scope / ``freezing'' limitation.} We agree that freezing mixing matrices is not standard end-to-end practice; it is a deliberate architectural choice that makes the deep mean-field analysis tractable while preserving expressive channelized dynamics. We will clarify that our target scope is theory-guided low-rank training (e.g.\ oscillatory/SciML settings), not a claim of covering all end-to-end deep learning.
  \item \textbf{Novelty vs.\ \NP{}.} We will make the novelty explicit: the contribution is not a new mean-field framework but the low-rank architecture + frozen mixing that avoids the deep i.i.d.\ collapse mechanism and yields a tractable $r$-channel ODE system, with a global-convergence conclusion within the parametrized class when the dynamics converges.
  \item \textbf{Theorem 4.1 and mechanism.} We agree the two-point toy is minimal; we will add a developed ``mechanism'' remark after Theorem~4.1 (feedback loop through $H_2=\sum_k L_{c_2,k}f_k$ and ReLU gates), and we will better separate what is proved from what is heuristic/empirical in higher dimension.
  \item \textbf{Symmetry sensitivity / hyperparameters.} We will add a short variance-over-seeds summary (and, where possible, a table/plot) clarifying which symmetry observable is measured and how it varies with $r$ and optimization settings. We will also adjust claims to avoid over-interpreting a conservative finite-width bound.
  \item \textbf{Exponential Grönwall factor.} The exponential prefactor in Theorem~3.3-type bounds is a worst-case Grönwall artifact; it should not be read as a predictor of practical symmetry breakdown. We will rephrase this part and avoid claiming that ``channel specialization avoids the exponential bound'' without evidence.
  \item \textbf{$(d,r,N)$ scaling.} We will add discussion and, resources permitting, synthetic sweeps varying intrinsic dimension and rank to probe whether/where a transition occurs.
\end{itemize}

\section*{Author response --- Reviewer Z9GA}

\textbf{Comment:}
\begin{verbatim}
As I mentioned in the Strengths and Weaknesses part, the architecture considered in this paper is a stacking of two-layer network, I don't see the difficulty of proving the mean-field limit when both
 are trained. Could the authors elaborate more on the technical difficulty?

I wonder how universal is spectral bias discovered in Section 4 for low-rank training? For example, if we do not consider network in the mean-field regime, do we still have such spectral bias?
\end{verbatim}
\textbf{Response:} We thank the reviewer for the clear summary of the technical concerns.
\begin{itemize}
  \item \textbf{Constant rank / one trained factor.} We agree this is not the most general practical setting; it is the setting in which we can (i) avoid deep i.i.d.\ collapse while (ii) keeping a closed mean-field description and a tractable global-convergence argument. We will clarify that this is a first step toward understanding low-rank optimization mechanisms in deep networks.
  \item \textbf{``Easier than \NP{}''.} Our architecture does simplify certain aspects (frozen mixing reduces degrees of freedom), but it introduces a new channelized coupling (the $r$-channel forward map $H_\ell=\sum_k L_{\cdot,k}f_k$ and corresponding backward signals), which is exactly where the proof differs: the Lipschitz/coupling chain carries $\|L^{(\ell)}\|_{\infty,1}$ and a $\max_k$ over channels. We will highlight the single step where the argument is not verbatim.
  \item \textbf{Training both factors.} Training both factors removes the frozen heterogeneity that breaks permutation symmetry and substantially complicates the coupling estimates (and may reintroduce collapse modes). We will state this obstacle explicitly and position it as open future work.
  \item \textbf{Spectral bias.} We agree Section~4 is currently qualitative. We will (i) tone down claims, (ii) clearly separate NTK/lazy-regime statements (where classical low-frequency bias persists) from the rich feature-learning regime where we observe different behavior, and (iii) add citations for the polyhedral/Fourier intuition (Rahaman et al.\ 2019; Mont\'ufar; Raghu) while labeling it as heuristic.
\end{itemize}

\section*{Author response --- Reviewer ByWV}

\textbf{Comment:}
\begin{verbatim}
Typo at (040) "...is low rank all we need for global convergence ? Low-rank networks...": There is an extra space between the question mark and the first word of the second sentence. Typo at (042): "...with
. substantially...".

At (044), it is posed that "can gradient-based training still converge to a global minimizer, or is full rank essential?". I didn't understand why improvement in the training performance, as shown in Fig.
, implies that the low-rank structure affects the dynamics of the gradient-based optimization algorithms such that they converge to a global minimizer. First, the low-rank factorization might change the loss landscape completely if some global minimizers are not feasible for the specific factorization determined by
. Note that the low-rank problem is different than the full-rank problem. Therefore, the functions represented by the global minimizers of the full-rank problem can be different than the low-rank problem.

For the Fig.
, how are different momentum rates selected for the corresponding level of rank constraints?

Do low-rank random feature networks differ from the multi-component multilayer neural networks proposed by Zhang et al. (2023) only in their low-rank structure?

Typo at (137),
 should be denoted by vector notation.

Could you clarify the convergence assumptions made by Nguyen & Pham (2023) for multilayer and two layer neural networks?

Nguyen & Pham, 2023 established the global convergence of multilayer neural networks trained under stochastic gradient descent (SGD). Is your convergence analysis also based on SGD?
\end{verbatim}
\textbf{Response:} We thank the reviewer for the concrete suggestions.
\begin{itemize}
  \item \textbf{Experiments.} We will expand the empirical section (additional datasets/sweeps where space permits) and add variance-over-seeds/frozen-feature draws to quantify sensitivity.
  \item \textbf{What is claimed about global minima.} We agree the current wording can be misleading. Our guarantee is \emph{within the low-rank parametrized class}: low-rank and full-rank problems have different feasible sets and need not share global minimizers. We will state this explicitly and avoid implying equality of full-rank and low-rank optima.
  \item \textbf{Momentum selection.} Chosen by a small sweep; we will report the sweep protocol.
  \item \textbf{Zhang et al.\ (2023).} The architecture matches multi-component networks; our novelty is the mean-field optimization analysis (channel ODEs, global-convergence conclusion within the parametrized class, and specialization mechanism). Zhang et al.\ do not analyze optimization.
  \item \textbf{Assumptions in main text.} We will move a compact assumption block (bounded activations/mixing, initialization, loss regularity, convergence-to-limit-point) into the main text with a short discussion.
  \item \textbf{What is new vs.\ \NP{}.} We will emphasize the single non-verbatim step: the low-rank channelized forward map introduces $\|L^{(\ell)}\|_{\infty,1}$ and $\max_k$ over channels in the coupling/Lipschitz chain, and frozen mixing breaks the deep i.i.d.\ collapse mechanism.
  \item \textbf{\NP{} convergence assumptions and SGD.} We will clarify that \NP{} use a weighted coupled-$L^1$ transport condition under a coupling $\pi_t$ (not merely moment convergence), and that our analysis is in the same mean-field/SGD scaling where SGD is a small-step discretization of the limiting ODE (gradient flow).
  \item \textbf{Typos / vector notation.} We will fix the listed typos and notation inconsistencies.
\end{itemize}

\section*{Author response --- Reviewer A6kZ}

\textbf{Comment:}
\begin{verbatim}
1. Convergence guarentees. Are there conditions under which convergence is guaranteed for RF-LR models (those mentioned in section 3.5 or otherwise)?

2. High-dimension extensions. Does the spike learning mechanism extend to higher dimensions with overlapping spikes in feature space?

3. Sensititivity of randomness of fixed feature maps. How sensitive are the convergence guarantees and empirical performance to the initial draw of the frozen feature maps for finite width networks?
\end{verbatim}
\textbf{Response:} We thank the reviewer for the thorough and constructive critique.
\begin{itemize}
  \item \textbf{``Channel specialization avoids the exponential bound''.} We agree the current wording is unsupported. We will remove or rephrase this as a conservative statement (the bound is worst-case), and we will add empirical evidence if we keep any specialization-related claim.
  \item \textbf{Minimum network size claim.} We will either add a citation/justification or remove the claim from Appendix~F.6.
  \item \textbf{Beyond two points/two channels.} We agree the toy theorem is limited. We will add a more developed treatment: a mechanism remark after Theorem~4.1 explaining the reinforcing-loop effect through $H_2=\sum_k L_{c_2,k}f_k$ and ReLU gates, and we will clarify that the ODE-level mechanism is dimension-agnostic (the toy proof is chosen for tractability).
  \item \textbf{Variance over frozen features.} We have not run a systematic multi-draw sweep over independent frozen-feature realizations; we will state this and, if space permits, add compact diagnostics or defer to future work.
  \item \textbf{Presentation fixes.} We will fix the ordering in Section~4.1 (present Theorem~4.1 before referencing it) and revise the figure captions (notably Figs.~8--9) for clarity.
  \item \textbf{Convergence guarantees.} Our main theorem is conditional: if the mean-field dynamics converges, the limit is a global minimizer within the parametrized class. Unconditional convergence rates/guarantees for RF-LR remain largely open; we will clarify this and point to the sharpest known partial results (and what is future work).
  \item \textbf{High-dimensional spike mechanism.} The channel feedback mechanism extends in the same form at the ODE level (same algebra of $H_2=\sum_k L_{c_2,k}f_k$ and channelwise backward signals), but the clean two-point proof does not directly transfer because high-dimensional features can overlap; we will clarify this distinction and add empirical support.
  \item \textbf{Sensitivity to frozen draw.} Guarantees use bounded frozen mixing and feature richness, not a specific law (finite-support draws suffice; Gaussians in code are convenience). Finite-width performance can vary with the draw; we have not isolated that variance empirically post-deadline. We will spell out theory vs.\ mean-field idealization in the camera-ready.
\end{itemize}

\end{document}

